\chapter{Client-Server-Struktur}

Das Ziel eines Messaging-Programmes ist es, von Gerät A eine Nachricht an Gerät B zu versenden. Dabei sollte möglichst wenig Rechenleistung beansprucht werden. Das Versenden soll schnell erfolgen, und die Nachricht darf für niemanden ausser den Absender und den Empfänger lesbar sein. Zum Versenden von Nachrichten gibt es schon sehr viele Programme. Die meisten können anhand ihres Funktionsprinzips in drei Unterkategorien eingeteilt werden: Diese sind die Client-Server-Struktur, das Peer-to-Peer-Messaging und das föderierte Servermodell. Die Client-Server-Struktur ist die verbreitetste Architektur für Messaging-Programme und auch diejenige, welche wir für unser Projekt benutzen. Deshalb werden wir in diesem Kapitel die Theorie dahinter erklären.
\section{Funktionsprinzip der Client-Server-Struktur}
Bekannte Beispiele für Programme, die die Client-Server-Struktur verwenden, sind WhatsApp, Signal und Discord. Dabei werden Clients benutzt, die sich alle mit einem Server verbinden \footnote{Dies kann auch einen Server Cluster }. Der Client ist ein Programm, welches auf dem Gerät jedes Nutzers läuft. Er dient dazu, Nachrichten an den Server zu senden und von ihm zu empfangen. Dazu hat er ein Interface, auf dem die Nachrichten angezeigt werden und neue verfasst werden können. Der Server ist eine zentrale Rechenmaschine, welche gleichzeitig mit allen Clients verbunden ist. Er speichert die vom Client verschickten Nachrichten ab, bis er eine Verbindung zum gewählten Empfänger hergestellt hat, und leitet sie dann an ihn weiter. Zusätzlich kann der Empfänger, wenn er die Nachricht empfängt und liest, dies dem Server mitteilen, und der Server leitet diese Information an den Sender zurück. Einige Programme, wie zum Beispiel Discord, löschen die Nachrichten nicht vom Server. Wenn dies gemacht wird kann der Nutzer alle vergangenen Nachrichten auf einem neuen Gerät abrufen.
\subsection{Der Aufbau}
\subsection{Verbindung zwischen Client und Server}
\subsection{Kryptologie}
Damit von einem Messaging-Programm Privatsphäre gewährleistet wird, darf niemand ausser Absender und Empfänger Zugriff auf die versendete Nachricht haben. Dazu wird ein Verfahren namens Verschlüsselung verwendet. Wenn Nachrichten ohne Verschlüsselung verschickt werden, erlaubt das jedem, der Zugriff auf das Netzwerk der beteiligten Clients oder des Servers hat, die Nachricht auch zu lesen.
\\\\
Die Definition von Wikipedia: 
\begin{quote}
	\textit{"Verschlüsselung nennt man die Umwandlung von Informationen, genannt Klartext, in einen Geheimtext.
     Dies geschieht mithilfe von kryptografischen Verfahren unter Verwendung eines geheimzuhaltenden Schlüssels,
     der nur den befugten Personen bekannt sein darf. Mit Kenntnis des Schlüssels kann der Geheimtext wieder in den Klartext zurückgewandelt werden,
     was Entschlüsselung genannt wird."
}
\end{quote}
Bei einer Ende-zu-Ende-Verschlüsselung wird garantiert, dass niemand ausser dem Absender und dem Empfänger die Nachricht lesen kann. Dabei wird die Nachricht am Ende des Absenders verschlüsselt und kann nur am Ende des Empfängers wieder entschlüsselt werden. Niemand dazwischen kann die Nachricht entschlüsseln. Der Server erhält also nur eine verschlüsselte Version der Nachricht. Um die verschlüsselte Nachricht wieder zu entschlüsseln, braucht der Empfänger jedoch zwingend den korrekten Schlüssel. Dies ist bei einer herkömmlichen symmetrischen Verschlüsselung, derselbe Schlüssel wie zur Verschlüsselung verwendet wurde. Somit entsteht das Problem, dass er dem Empfänger mitgeteilt werden muss. Durch das Versenden des Schlüssels hätte jedoch jeder Zugriff auf den Schlüssel und die verschlüsselte Nachricht. Die Nachricht könnte dann von allen entschlüsselt werden. Um dieses Problem zu lösen, wird eine asymmetrische Verschlüsselung benutzt.
\subsubsection{asymmetrische Verschlüsselung}
Die asymmetrische Verschlüsselung zeichnet sich dadurch aus, dass sie nicht den gleichen Schlüssel zur Verschlüsselung und Entschlüsselung verwendet. Dafür gibt es verschiedene Algorithmen, die alle dem gleichen Prinzip folgen. Bei der Erstellung eines Accounts generiert der Client A ein sogenanntes Key-Pair. Dieses besteht aus einem Public- und einem Private-Key. Der Public Key wird offen an alle anderen Nutzer versendet, der Private Key wird lediglich auf dem Gerät des Nutzers A gespeichert. Wenn Client B eine Nachricht an Client A versenden möchte, benutzt er den Public Key von Client A, um diese zu verschlüsseln. Der Public Key kann aber nicht verwendet werden, um die Nachricht wieder zu entschlüsseln. Eine Entschlüsselung ist nur mit dem Private Key von Client A möglich. So wird eine sichere Kommunikation ermöglicht, ohne die Schwachstelle, dass ein Schlüssel zur Entschlüsselung ausgetauscht werden muss. Solange alle Private Keys geheim bleiben, ist völlige Privatsphäre gewährleistet. Der Nachteil an diesem Verfahren ist, dass die Algorithmen für die asymmetrische Verschlüsselung viel langsamer sind und mehr Rechenleistung beanspruchen.
\subsubsection{Hybrid-Kryptosystem}
Die wichtigste Art der Verschlüsselung ist das Hybrid-Kryptosystem. Es ist, wie der Name schon andeutet, eine Kombination aus der asymmetrischen und der symmetrischen Verschlüsselung. Während dem Herstellen einer Verbindung zwischen Client A und Server wird von Client A ein zufälliger Session Key generiert. Dieser wird dann benutzt, um die Nachricht, welche den Grossteil der Datenmenge ausmacht, effizient symmetrisch zu verschlüsseln. Danach wird der Session Key asymmetrisch verschlüsselt mit dem Public Key von Client B. Da der Session Key sehr klein ist, wird trotz asymmetrischer Verschlüsselung wenig Zeit und Rechenleistung beansprucht. Danach werden die verschlüsselte Nachricht und der verschlüsselte Session Key an Client B versendet. Dieser entschlüsselt zuerst den Session Key mit seinem Private Key und kann dann die Nachricht mit dem Session Key symmetrisch entschlüsseln. Bei diesem Vorgang wird die grosse Datenmenge effizient verschlüsselt mit der symmetrischen Verschlüsselung und ist trotzdem sicher, da der Session Key asymmetrisch verschlüsselt wird. Das Hybrid-Kryptosystem ist der heutige Standard und wird unter anderem für HTTPS gebraucht, das Protokoll das hauptsächlich genutzt wird, um im Internet Daten zu übertragen.
\section{Vergleich mit anderen Strukturen}
